% Modernised reading — fonds Alexandre Grothendieck, shelfmark 55, pages 1-17.
%
% This is an INTERPRETATION, not a transcription. Everything the critical
% apparatus carried moves into footnotes. In any dispute of substance the
% transcription governs: whoever cites Grothendieck must cite batch-01.fr.tex.
%
% Pass: Opus 5 (claude-opus-5), 2026-09-13 — first pass, unchecked against the
% pages by a human. The folder was read whole; it holds one batch, pages 1-17,
% of which page 15 carries nothing.
%
% What the folder is. Three things bound together. (1) Manuscript notes, pages
% 1, 3, 5, 7, 9, 11, 13: over a Stein compact K, with A = Gamma(K, O), the
% functor M |-> M~ is an equivalence from finitely presented A-modules to
% coherent sheaves; abstracted to a ringed space (Lemme A), extended to
% algebras of finite presentation, and closed by a noetherianity statement
% (Lemme B) whose corollary — A noetherian for a compact polycylinder — is
% Frisch's theorem. (2) Six annotated typescript leaves on dimension theory
% (pages 2, 4, 6, 10, 12, 14) and one on associated primes (page 8). (3) Pages
% 16-17: a two-dimensional complex torus fibred in elliptic curves over an
% elliptic curve which is not an abelian variety.
%
% Notation fixed for the whole document. K a Stein compact, O = O_K the sheaf
% of germs of holomorphic functions along K; for a ringed space X, O_X and
% A = Gamma(X, O_X); M~ = M (x)_{A_X} O_X, A_X the constant sheaf. B (script)
% an O_X-algebra, B_m its sections of degree <= m when B = O_X[T_1..T_n]. He
% underlines the O; written plainly here.
%
% Substantive departures from the pages, each footnoted where it occurs:
% — page 1, « X espace de Stein compact »: a compact Stein space is a finite
%   set; what the notes use is a Stein compact (a compact set with a basis of
%   Stein neighbourhoods) with the sheaf of germs. The reading says so.
% — page 1, « modules de présentation finie sur X »: read « sur A ».
% — page 5, Lemme A: under (i)-(iii) as written the proof needs global
%   generation of coherent sheaves (used on page 3 as « on sait »), which an
%   abstract ringed space does not provide; the reading states the lemma with
%   (beta) added, which is what page 5's own Remarque does.
% — page 11: passing Gamma to the limit needs Gamma to commute with filtered
%   colimits (true for K compact Hausdorff, or X noetherian); supplied.
% — pages 16-17: the lattice written, generated by e1, ie1, e2, e1 + ik e2,
%   contains ik e2, so T = C/Z[i] x C/(Z + ikZ) IS an abelian variety, and the
%   argument fails exactly when x lies on the e2-axis (where « c', c''
%   rationnels » does not follow). The claim is true; the reading gives the
%   lattice e1, ie1, e2, k e1 + i e2, for which the page's computation goes
%   through, with the roles of the two coordinates exchanged.
% — page 16: « V soit Gamma-propre » where V' is meant.
% — page 4 (typescript margin): « dim O_x <= dim O_y + dim O_x » read as the
%   dimension inequality dim O_x <= dim O_y + dim O_{X_y, x}.
% — page 14 (typescript): replacing « équicodimensionnel » by
%   « équidimensionnel » makes d) false (counterexample Spec R[x], R a DVR),
%   and no form of b) is equivalent to a), not even with both notions and the
%   chain condition (a six-point finite space refutes it). The reading keeps
%   a), c), d) in their true form and says why b) is dropped.
%
% Modern names given and footnoted as ours: compact de Stein, théorèmes A et
% B de Cartan, anneau cohérent, théorème de Frisch (and Siu's criterion),
% théorème de réductibilité de Poincaré, tores de dimension algébrique 1,
% variété de Hodge / Kodaira, espace caténaire, équicodimensionnel.
%
% What the folder announces and never establishes: Corollaire 3 of page 5
% (illegible); the coherence of the J_m (« ? histoire d'anneaux cohérents ? »)
% and with it (beta) for algebras of finite presentation; the proof of Lemme
% B; the answer to page 7's question beyond « ok. dans le cas classique ».
\documentclass[11pt,a4paper]{article}
% Le préambule commun à toutes les transcriptions du fonds Grothendieck.
%
% Il tient en une idée : **l'incertitude se compose, elle ne se résout pas.**
% Une transcription qui lisse un mot douteux a détruit la seule chose pour
% laquelle on la faisait — un lecteur doit pouvoir savoir, à chaque endroit,
% ce qui était sur le feuillet et ce qui a été ajouté. Les six macros qui
% suivent sont donc l'appareil critique, et non de la décoration.
%
% Les mêmes commandes sont reconnues par `scripts/render.mjs`, qui produit la
% vue de lecture du site. Une macro ajoutée ici doit l'être aussi là-bas,
% faute de quoi le rendu échouera bruyamment — ce qui est voulu : un rendu
% silencieusement faux est pire qu'un rendu absent.

\ProvidesPackage{grothendieck}[2026/08/08 Transcriptions du fonds Grothendieck]

\usepackage{amsmath,amssymb,amsthm}
\usepackage{mathrsfs}
\usepackage{stmaryrd}
% Les diagrammes commutatifs sont partout dans ce fonds : c'est la langue
% même de la géométrie algébrique de Grothendieck, et une transcription qui
% les rendrait en prose ne transcrirait rien.
\usepackage{tikz-cd}
% Français par défaut — c'est la langue des feuillets — mais l'anglais est
% chargé aussi : la traduction et le résumé sont des documents anglais, et la
% typographie française (espaces avant les deux-points) y serait une faute.
% Ces documents basculent par \selectlanguage{english} en tête de corps.
\usepackage[english,french]{babel}
\usepackage{fontspec}
\usepackage{geometry}
\geometry{a4paper,margin=25mm,marginparwidth=28mm}
\usepackage{marginnote}
\usepackage{xcolor}
% ulem, not soul. `soul`'s \st and \ul are built on a character-by-character
% scan that XeLaTeX's font machinery breaks: the document fails with an
% undefined control sequence rather than degrading. ulem does the same two jobs
% and is Unicode-engine safe. `normalem` stops it hijacking \emph, which the
% transcriptions use for Grothendieck's own emphasis.
\usepackage[normalem]{ulem}
\usepackage{hyperref}
\hypersetup{colorlinks=true,linkcolor=black,urlcolor=black,pdfborder={0 0 0}}

% --- Métadonnées du lot -----------------------------------------------------
% Reprises telles quelles par le script de rendu, qui les affiche en tête de
% la vue de lecture. Elles fixent ce que le fichier prétend couvrir : sans
% elles, une transcription flotte sans point d'attache dans un fonds de seize
% mille feuillets.

\newcommand{\folder}[1]{\gdef\grothFolder{#1}}
\newcommand{\batch}[1]{\gdef\grothBatch{#1}}
\newcommand{\leaves}[2]{\gdef\grothFirst{#1}\gdef\grothLast{#2}}
\newcommand{\foldertitle}[1]{\gdef\grothFoldertitle{#1}}
\gdef\grothFolder{?}\gdef\grothBatch{?}\gdef\grothFirst{?}\gdef\grothLast{?}\gdef\grothFoldertitle{}

% --- Le résumé --------------------------------------------------------------
%
% Il ouvre la lecture modernisée, sous le titre « Résumé », et il est composé
% en petit. Ce n'est pas un ornement : c'est le seul passage du document qui
% ne soit pas des mathématiques, et un lecteur doit pouvoir le reconnaître
% avant de l'avoir lu — puis le sauter s'il connaît déjà le sujet, ou s'y
% arrêter s'il ne le connaît pas. Le filet qui le suit marque la reprise du
% corps.

\newenvironment{resume}
  {\small\setlength{\parskip}{0.45em}}
  {\par\medskip\hrule\medskip}

% --- Mots-clés ---------------------------------------------------------------
%
% En anglais, en clôture du résumé de la lecture modernisée : le vocabulaire
% moderne sous lequel chercher ce que les feuillets construisent. Le manifeste
% du site (`npm run manifest`) les extrait pour étiqueter la cote entière —
% cette ligne est la source unique des tags, il n'y a pas de fichier à côté.
\newcommand{\keywords}[1]{\par\smallskip\noindent
  {\footnotesize\textsc{Keywords} --- \textit{#1}}\par}

% --- L'appareil critique ----------------------------------------------------

% Le numéro de feuillet, en marge. C'est l'ancre qui permet de régler un
% désaccord sur une formule en regardant *un* feuillet plutôt qu'une chemise
% de six cents. Le site s'en sert aussi pour tourner les pages du fac-similé
% quand on fait défiler la transcription.
\newcommand{\leaf}[1]{%
  \marginnote{\footnotesize\color{gray!70}\textsf{#1}}%
  \label{leaf:#1}\ignorespaces}

% Illisible : on ne devine pas. Un mot inventé qui se lit comme les autres est
% le pire résultat possible.
\newcommand{\ill}{\textcolor{red!60!black}{[\,\ldots\,]}}

% Lecture douteuse : le mot est donné, et signalé comme incertain.
\newcommand{\uncertain}[1]{\uline{#1}}

% Ajout de l'éditeur — une lettre manquante, un mot sous-entendu.
\newcommand{\add}[1]{\textcolor{blue!50!black}{[#1]}}

% Biffé par Grothendieck lui-même. Ce qu'il a rayé fait partie du document :
% une rature dit souvent où la pensée a changé de direction.
\newcommand{\struck}[1]{\sout{#1}}

% Note du transcripteur, sur l'état matériel du feuillet.
\newcommand{\note}[1]{\par\smallskip{\small\color{gray!60!black}\textit{[#1]}}\par\smallskip}

% Note marginale de Grothendieck, distincte du corps du texte.
\newcommand{\marginal}[1]{\par\smallskip\noindent\hspace{1em}%
  {\small\color{gray!60!black}#1}\par\smallskip}

% --- Titre ------------------------------------------------------------------

\AtBeginDocument{%
  \begin{center}
    {\large\bfseries Fonds Alexandre Grothendieck --- cote n\textsuperscript{o}~\grothFolder}\\[2pt]
    {\itshape\grothFoldertitle}\\[4pt]
    {\small feuillets \grothFirst--\grothLast\ (lot \grothBatch)}\\[2pt]
    {\footnotesize\color{gray!60!black}Transcription établie d'après le fac-similé numérisé
     par l'Université de Montpellier.\\
     Les lectures incertaines sont soulignées, les illisibles notées [\,\ldots\,],
     les ajouts entre crochets.}
  \end{center}
  \vspace{1em}}

\endinput


\folder{55}
\pages{1}{17}
\dating{[à partir de 1961]}
\watermark{Édition de démonstration\\interprétation personnelle de l'œuvre}
\foldertitle{Espaces analytiques : notes manuscrites (s.d.) — lecture modernisée du dossier entier}

\begin{document}

\section*{Résumé}

\begin{resume}
En géométrie algébrique, un espace affine se lit tout entier sur l'anneau de
ses fonctions : les objets géométriques qui vivent dessus --- les faisceaux
cohérents --- ne sont rien d'autre que des modules sur cet anneau, et l'on
peut passer d'un langage à l'autre sans rien perdre. Ce dictionnaire est la
raison d'être de la notion de schéma affine. La question que posent ces notes
est de savoir s'il existe un dictionnaire semblable en géométrie analytique
complexe, où les fonctions sont holomorphes et non polynomiales.

Sur un ouvert de $\mathbb{C}^n$ la réponse est non : l'anneau des fonctions
holomorphes est trop gros, et un faisceau cohérent peut demander une infinité
de sections pour être engendré. Mais sur un \emph{compact} convenable --- un
polydisque fermé, avec les germes de fonctions holomorphes au voisinage ---
les théorèmes de Cartan donnent exactement ce qu'il faut : tout faisceau
cohérent est engendré par un nombre fini de sections globales, et la
cohomologie s'annule. Les notes montrent qu'alors le dictionnaire marche
(pages 1 et 3), puis isolent ce qui le fait marcher, pour un espace annelé
quelconque (page 5), puis l'étendent aux algèbres (pages 7, 9 et 11). Elles se
ferment sur la question qui décide de tout : l'anneau $A$ des fonctions sur le
compact est-il noethérien, comme l'est celui d'une variété affine ? Un lemme
le démontre sous des hypothèses qu'il écrit, et la marge note « premier
espoir ». Pour un polydisque fermé, c'est un théorème de Frisch, publié
quelques années plus tard.

Le dossier contient aussi deux choses sans rapport avec ce qui précède. Des
feuillets dactylographiés d'un traité sur la dimension, qu'il a annotés en
correcteur exigeant (« propositions laborieuses inutiles », « abusif »). Et un
petit fragment, pages 16 et 17, qui donne un exemple : un tore complexe de
dimension 2, construit comme le plan $\mathbb{C}^{2}$ enroulé sur un réseau, qui est un
fibré en courbes elliptiques au-dessus d'une courbe elliptique, et qui
pourtant ne peut être plongée dans aucun espace projectif. L'idée tient à un
nombre irrationnel glissé dans le réseau. Le réseau que la page écrit ne
convient pas --- il se scinde, et le tore obtenu est un produit --- mais une
permutation du nombre irrationnel suffit à le réparer, et le calcul de la page
passe alors tel quel.

\keywords{Stein compact, coherent sheaf, Cartan theorems A and B, coherent ring, Frisch theorem, non-algebraic complex torus, Poincaré reducibility, Krull dimension, catenary space}
\end{resume}

\section*{Le fil du dossier, et les conventions}

\emph{Les notes manuscrites}, à l'encre bleue, alternent feuillet à feuillet
avec les tapuscrits et suivent un seul fil :

\begin{enumerate}
\item[1.] \emph{Le dictionnaire sur un compact de Stein} (pages 1 et 3).
\item[2.] \emph{Le lemme A} (page 5) : ce qu'il faut à un espace annelé pour
que le dictionnaire tienne ; ses corollaires, exactitude, platitude,
cohérence de $A$.
\item[3.] \emph{Les algèbres} (page 7), et la question de pure algèbre
commutative qui en sort.
\item[4.] \emph{Les algèbres de polynômes} (page 11, puis page 9) : passage à
la limite sur les degrés. L'ordre des archivistes n'est pas ici celui de la
rédaction, et on lit la page 11 avant la page 9.
\item[5.] \emph{Le lemme B} (page 13) : noethérianité de $A$.
\end{enumerate}

\emph{Les tapuscrits} (pages 2, 4, 6, 10, 12, 14 et 8) sont lus à part, dans
l'ordre de leur propre pagination, pour ce qu'il y a ajouté. \emph{Le tore}
(pages 16 et 17) est lu en dernier.

\emph{Conventions.} Pour un espace annelé $X$, on note $\mathcal{O}_{X}$ son
faisceau structural, $A = \Gamma(X, \mathcal{O}_{X})$, $A_{X}$ le faisceau
constant d'anneaux, et pour un $A$-module $M$,
$\widetilde{M} = M \otimes_{A_{X}} \mathcal{O}_{X}$. Un \emph{compact de
Stein} est un compact $K$ d'un espace analytique complexe qui admet un système
fondamental de voisinages de Stein ; on le munit du faisceau
$\mathcal{O}_{K}$ des germes de fonctions holomorphes le long de $K$.
$\mathcal{B}$ est une $\mathcal{O}_{X}$-algèbre, et, pour
$\mathcal{B} = \mathcal{O}_{X}[T_{1}, \ldots, T_{n}]$, $\mathcal{B}_{m}$ est
le sous-Module des polynômes de degré au plus $m$, libre de rang fini sur
$\mathcal{O}_{X}$.

\emph{Ce que le dossier annonce et n'établit pas} : le corollaire 3 de la page
5, sur les idéaux maximaux, presque illisible ; la cohérence des idéaux
tronqués $J_{m}$ de la page 9, et avec elle l'engendrement par les sections
pour les algèbres de présentation finie ; la démonstration du lemme B ; la
réponse à la question de la page 7 au-delà du cas « classique ».

\pagerange{1}{3}
\section*{Le dictionnaire sur un compact de Stein (pages 1 et 3)}

Soit $K$ un compact de Stein et $A = \Gamma(K, \mathcal{O}_{K})$.\footnote{La
page écrit « $X$ espace de Stein compact ». Un espace de Stein compact est
un ensemble fini de points, et rien dans la suite n'aurait d'intérêt ; ce que
les notes utilisent --- l'annulation de la cohomologie cohérente et
l'engendrement par les sections globales, sur un objet compact --- est la
situation d'un compact admettant une base de voisinages de Stein, avec le
faisceau des germes le long du compact. C'est aussi ce que confirme la page 13,
où « espace de Stein » est biffé au profit de « polycylindre analytique
(compact) ». On lit ainsi. Une première ligne, « $B \in$ Alg.\ de présentation
finie », est biffée en tête de page ; elle annonce la page 7.}

\textbf{Énoncé.} Le foncteur $M \mapsto \widetilde{M}$ est une équivalence de
la catégorie des $A$-modules de présentation finie sur celle des
$\mathcal{O}_{K}$-Modules cohérents.\footnote{La page écrit « modules de
présentation finie sur $X$ » ; le sens demande $A$. Un premier départ,
encadré puis biffé, voulait prouver la fidélité par
« $\widetilde{M} = 0 \Rightarrow M = 0$ » à partir d'une exactitude suivie
d'un point d'interrogation ; il est abandonné pour l'argument qui suit.}

Deux faits sur $K$ servent, et ce sont les théorèmes de Cartan :
\begin{enumerate}
\item[(A)] tout $\mathcal{O}_{K}$-Module cohérent $F$ est engendré par un
nombre fini de sections globales : il existe une surjection
$\mathcal{O}_{K}^{q} \to F$ ;
\item[(B)] $H^{1}(K, F) = 0$ pour tout $F$ cohérent.
\end{enumerate}

\emph{$M \simeq \Gamma(\widetilde{M})$.} Une présentation
$A^{p} \to A^{q} \to M \to 0$ donne, le produit tensoriel étant exact à
droite, $\mathcal{O}^{p} \to \mathcal{O}^{q} \to \widetilde{M} \to 0$. Par
(B), appliqué au noyau et à l'image, qui sont cohérents, la suite des sections
reste exacte ; comparant les deux suites, $M \xrightarrow{\sim}
\Gamma(\widetilde{M})$.

\emph{Pleine fidélité} (page 3). Pour $N$ un autre module de présentation
finie, on a le diagramme commutatif à lignes exactes

\begin{tikzcd}
0 \arrow[r] & \mathrm{Hom}(M, N) \arrow[r] \arrow[d, "\alpha"] & N^{q} \arrow[r] \arrow[d, "\beta"] & N^{p} \arrow[d, "\gamma"] \\
0 \arrow[r] & \mathrm{Hom}(\widetilde{M}, \widetilde{N}) \arrow[r] & \Gamma(\widetilde{N})^{q} \arrow[r] & \Gamma(\widetilde{N})^{p}
\end{tikzcd}

obtenu en appliquant $\mathrm{Hom}(-, N)$ et $\mathrm{Hom}(-, \widetilde{N})$
aux deux présentations. $\beta$ et $\gamma$ sont bijectives par ce qui
précède, donc $\alpha$ l'est aussi.

\emph{Surjectivité essentielle.} Si $F$ est cohérent, (A) donne
$\mathcal{O}^{q} \to F \to 0$ ; le noyau $R$ est cohérent, et (A) encore donne
$\mathcal{O}^{p} \to R \to 0$, d'où une présentation
\[
  \mathcal{O}^{p} \xrightarrow{u} \mathcal{O}^{q} \to F \to 0 .
\]
La matrice $u$ est faite de sections globales, donc définit
$u : A^{p} \to A^{q}$, et pour $M = \mathrm{Coker}\, u$ on a
$F \simeq \widetilde{M}$.

\pagerange{5}{5}
\section*{Le lemme A (page 5)}

La page 5 abstrait ce qui vient d'être fait.

\textbf{Lemme A.} Soit $X$ un espace annelé tel que
\begin{enumerate}
\item[(i)] $X$ est quasi-compact ;
\item[(ii)] $\mathcal{O}_{X}$ est un faisceau d'anneaux cohérent ;
\item[($\alpha$)] pour toute suite exacte $\mathcal{O}_{X}^{p} \to
\mathcal{O}_{X}^{q} \to F \to 0$, la suite des sections globales est exacte ;
\item[($\beta$)] tout $\mathcal{O}_{X}$-Module de présentation finie est
engendré par ses sections globales.
\end{enumerate}
Alors, avec $A = \Gamma(X, \mathcal{O}_{X})$, le foncteur
$M \mapsto \widetilde{M}$ est une équivalence de la catégorie des $A$-modules
de présentation finie sur celle des $\mathcal{O}_{X}$-Modules de présentation
finie, qui sont les cohérents.

La démonstration est celle des pages 1 et 3, mot pour mot : ($\alpha$) donne
$M \simeq \Gamma(\widetilde{M})$ et la pleine fidélité, ($\beta$) et (i) la
surjectivité essentielle.\footnote{La page énonce le lemme sous (i), (ii)
« $\mathcal{O}_{X}$ cohérent, les $\mathcal{O}_{X,x}$ noethériens », et (iii)
« $H^{1}(X, F) = 0$ pour tout $F$ cohérent ». (ii) et (iii) entraînent bien
($\alpha$), en appliquant (iii) à l'image et au noyau de $u$, cohérents par
(ii). Mais la surjectivité essentielle, à la page 3, invoque l'engendrement
par les sections comme un fait connu --- c'est le théorème A dans le cas de
Stein ---, et rien dans (i)--(iii) ne le fournit pour un espace annelé
quelconque : dans le cas analytique on le déduit de (B) en utilisant l'Idéal
cohérent d'un point, dont un espace annelé abstrait ne dispose pas. La
\emph{Remarque} qui clôt la même page propose précisément de remplacer (ii) et
(iii) par ($\alpha$) et ($\beta$) ; on énonce le lemme sous cette forme, en
gardant la cohérence de $\mathcal{O}_{X}$ sans laquelle « présentation finie »
et « cohérent » ne coïncident pas. La noethérianité des anneaux locaux ne sert
pas ici ; elle servira au lemme B.} Une \emph{variante}, écrite de bas en haut
dans la marge gauche, remplace la quasi-compacité par un système fondamental
de voisinages quasi-compacts et vise des présentations
$\mathcal{O}^{I} \to \mathcal{O}^{J}$ quelconques ; elle n'est pas achevée.

\textbf{Corollaire 1.} Le foncteur $M \mapsto \widetilde{M}$ est exact. En
particulier, pour tout $x \in X$, $\mathcal{O}_{X,x}$ est plat sur $A$.

\textbf{Corollaire 2.} L'anneau $A$ est cohérent : pour tout homomorphisme
$A^{n} \to A$, le noyau est de type fini.

On démontre le second d'abord. Le noyau $\mathcal{K}$ de
$\mathcal{O}_{X}^{n} \to \mathcal{O}_{X}$ est cohérent, donc de la forme
$\widetilde{N}$ avec $N$ de présentation finie, et
$N = \Gamma(\mathcal{K}) = \mathrm{Ker}(A^{n} \to A)$, $\Gamma$ étant exact à
gauche. Pour le premier : si $M' \to M$ est injectif entre modules de
présentation finie, le noyau de $\widetilde{M'} \to \widetilde{M}$ est un
$\widetilde{N}$, et $N = \Gamma(\widetilde{N}) = \mathrm{Ker}(M' \to M) = 0$.
Comme $\widetilde{M}_{x} = M \otimes_{A} \mathcal{O}_{X,x}$, la tensorisation
par $\mathcal{O}_{X,x}$ conserve l'injectivité $I \to A$ pour tout idéal de
type fini $I$ --- de présentation finie puisque $A$ est cohérent ---, et
$\mathcal{O}_{X,x}$ est plat.\footnote{Les deux démonstrations sont les
nôtres ; la page énonce les corollaires sans preuve. L'ordre inverse est
nécessaire : l'exactitude à gauche du foncteur, prise dans la catégorie de
tous les $A$-modules, demande que les noyaux calculés là soient de présentation
finie, c'est-à-dire que $A$ soit cohérent.}

\textbf{Corollaire 3.} Il porte sur les idéaux maximaux de type fini de $A$ et
les idéaux maximaux des $\mathcal{O}_{X,x}$, sous une hypothèse que $X$ soit
« pur ».\footnote{L'énoncé est presque illisible, et le mot « pur » n'est pas
assuré. On y lit « idéaux maximaux », « de type fini », « maximal de
$\mathcal{O}_{X,x}$ », et, dans une ligne serrée et en partie biffée,
« $(x, \mathfrak{m})$ », « $\mapsto x$ » ; le sens paraît être une
correspondance entre points de $X$ et idéaux maximaux de type fini de $A$,
mais cette lecture ne l'affirme pas.}

\pagerange{7}{7}
\section*{Les algèbres, et une question d'algèbre commutative (page 7)}

Du lemme A la page tire, pour les algèbres, deux conclusions :
\begin{itemize}
\item $B \mapsto B \otimes_{A_{X}} \mathcal{O}_{X}$ est une équivalence de la
catégorie des $A$-algèbres de présentation finie sur celle des
$\mathcal{O}_{X}$-Algèbres de présentation « globale » finie ;
\item pour une telle $B$, d'Algèbre associée $\mathcal{B}$,
$M \mapsto M \otimes_{B_{X}} \mathcal{B}$ est une équivalence des $B$-modules
de présentation finie sur les $\mathcal{B}$-Modules de présentation
« globale » finie.
\end{itemize}
Une $\mathcal{O}_{X}$-Algèbre est de présentation globale finie quand elle est
quotient de $\mathcal{O}_{X}[T_{1}, \ldots, T_{n}]$ par un Idéal engendré par
un nombre fini de sections globales.\footnote{L'adjectif « globale » est lu
avec doute, deux fois ; la définition est la nôtre, et c'est la seule qui
rende les deux énoncés vrais sous le seul lemme A : une présentation par des
sections globales est exactement ce qui provient d'une présentation sur $A$.
La page dit tirer ces conclusions « sans utiliser l'hyp.\ ($\alpha$) », ce
qu'on ne voit pas, la pleine fidélité en dépendant.}

\emph{Question.} Pour une $\mathcal{O}_{X}$-Algèbre $\mathcal{B}$ de
présentation finie, le noyau $\mathcal{J}$ de
$\mathcal{O}_{X}[T_{1}, \ldots, T_{n}] \to \mathcal{B}$ est-il cohérent ? Il y
voit une question « de pure algèbre commutative » : $A$ étant un anneau
cohérent et $J$ un idéal de type fini de $B = A[T]$, les
$J \cap B_{m}$ sont-ils des $A$-modules de type fini ? et, plus généralement,
les sous-$A$-modules de type fini de $B$ ont-ils cette
propriété ?\footnote{Si $A$ est noethérien, la réponse est oui sans rien
faire, $B_{m}$ étant un $A$-module libre de rang fini. Pour $A$ seulement
cohérent, elle n'est pas oui en général : un anneau de polynômes sur un anneau
cohérent n'est pas nécessairement cohérent (J.-P.~Soublin, 1970). On ne
précise pas ici le rapport exact entre la question de la page et la cohérence
de $A[T]$.} Une note au pied de la page répond pour un cas : « ok.\ dans le cas
“classique” $X$ polycylindre analytique ».

\pagerange{11}{9}
\section*{Les algèbres de polynômes (page 11, puis page 9)}

\textbf{Corollaire.} Supposons que $X$ satisfasse, en plus de (i),
\begin{enumerate}
\item[(ii')] les $\mathcal{O}_{X,x}$ sont noethériens, et
$\mathcal{O}_{X}[T_{1}, \ldots, T_{n}]$ est un faisceau d'anneaux cohérent
pour tout $n$ ;
\item[(iii')] tout point a un système fondamental de voisinages quasi-compacts
qui satisfont la condition (iii) de la page 5.
\end{enumerate}
Alors la conclusion du lemme A reste valable pour
$\mathcal{B} = \mathcal{O}_{X}[T_{1}, \ldots, T_{n}]$, et plus généralement
pour toute Algèbre de présentation finie.\footnote{Le titre est écrit trois
fois, « Corollaire », « Proposition », « Corollaire » ; un grand point
d'interrogation est porté en marge de (ii').}

\emph{La condition ($\alpha$) pour $\mathcal{B}$.} Soit
\[
  \mathcal{B}^{p} \xrightarrow{u} \mathcal{B}^{q} \to \mathcal{M} \to 0
\]
une suite exacte, et $N$ un majorant des degrés des coefficients $u_{ij}$. Elle
est limite inductive des suites exactes de $\mathcal{O}_{X}$-Modules
cohérents
\[
  \mathcal{B}_{m}^{p} \to \mathcal{B}_{m+N}^{q} \to \mathcal{M}_{m} \to 0 ,
\]
$\mathcal{M}_{m}$ désignant le conoyau. Chacune reste exacte sur les sections
par ($\alpha$) pour $\mathcal{O}_{X}$, et en passant à la limite
\[
  \Gamma(\mathcal{B}^{p}) \to \Gamma(\mathcal{B}^{q}) \to \Gamma(\mathcal{M}) \to 0
\]
est exacte : c'est ($\alpha$) pour $\mathcal{B}$.\footnote{Le passage à la
limite demande que $\Gamma$ commute aux limites inductives filtrantes, ce qui
est vrai sur un espace compact séparé --- le cas d'un compact de Stein --- ou
sur un espace noéthérien, mais pas sur un espace quasi-compact quelconque ;
la page ne le dit pas, et on l'ajoute.}

La démonstration continue en haut de la page 9.\footnote{La page 11 s'arrête
sur « Ceci prouve », et la page 9 commence par « … que le foncteur
$M \mapsto \widetilde{M}$ … est pl.\ fidèle … (en utilisant seulement la
condition $\alpha$) » ; on lit la seconde comme la suite de la première.
L'ordre des archivistes n'est pas celui de la rédaction ; on ne le corrige pas
dans la transcription, on le suit ici.} ($\alpha$) donne, comme au lemme A, la
pleine fidélité de $M \mapsto \widetilde{M}$ des $\Gamma(\mathcal{B})$-modules
de présentation finie vers les $\mathcal{B}$-Modules de présentation finie.
Reste ($\beta$) : qu'un $\mathcal{B}$-Module de présentation finie soit
engendré par ses sections. On voudrait l'obtenir des Modules cohérents
$\mathcal{M}_{m}$ ; mais l'image de $\mathcal{B}_{m+N}^{q}$ dans $\mathcal{M}$
est le quotient de $\mathcal{B}_{m+N}^{q}$ par son intersection avec l'image
de $u$, et la cohérence de ces intersections est la question de la page 7.
Écrivant $\mathcal{J} = \varinjlim_{m} \mathcal{J}_{m}$, il faut prouver que
les $\mathcal{J}_{m} = \mathcal{J} \cap \mathcal{B}_{m}$ sont
cohérents.\footnote{L'explication de pourquoi la question revient est la
nôtre ; la page écrit la limite et ce qu'il faut prouver, avec de nombreux
mots illisibles, et note en marge « ? histoire d'anneaux cohérents ? ». Un
N.B.\ encadré, qui supposait $X$ « pur », est annulé.} Le dossier ne le
prouve pas.

\pagerange{13}{13}
\section*{Le lemme B : noethérianité (page 13)}

\textbf{Lemme B.} Supposons que $X$ satisfasse (i), (ii), et
\begin{enumerate}
\item[(iv)] les anneaux locaux $\mathcal{O}_{X,x}$ sont noethériens ;
\item[(v)] \emph{principe d'identité} : pour tout Idéal cohérent $\mathcal{J}$,
posant $Y = (\mathrm{Supp}\, \mathcal{O}_{X}/\mathcal{J},
\mathcal{O}_{X}/\mathcal{J})$, et pour tout $y \in Y$ tel que
$\mathcal{O}_{Y,y}$ soit intègre, il existe un voisinage ouvert $U$ de $y$
dans $Y$ tel que toute section de $\mathcal{O}_{Y}$ sur $U$ nulle au voisinage
de $y$ soit nulle sur $U$.
\end{enumerate}
Alors, pour tout $F$ cohérent, toute suite croissante de sous-Modules cohérents
de $F$ est stationnaire.

\textbf{Corollaire.} Sous (i) à (v), et sous les hypothèses du lemme A,
$A = \Gamma(X, \mathcal{O}_{X})$ est noethérien.

Une suite croissante d'idéaux de type fini $I_{k}$ de $A$ donne, par le
Corollaire 1 de la page 5, une suite croissante de sous-Idéaux cohérents
$\widetilde{I_{k}}$ de $\mathcal{O}_{X}$, stationnaire par le lemme B ; et
$I_{k} = \Gamma(\widetilde{I_{k}})$. La condition des chaînes ascendantes pour
les idéaux de type fini suffit.\footnote{Le passage du lemme au corollaire est
le nôtre, et il utilise le lemme A, que la page ne cite pas mais dont la
notation $A$ provient. La démonstration du lemme B n'est pas sur la page ; la
condition (v) est ce qui, dans le cas analytique, permet de passer de la
stabilisation des germes en un point à la stabilisation sur un voisinage, par
la décomposition en composantes irréductibles et le principe du prolongement
analytique. Cette lecture ne vérifie pas que (i)--(v) suffisent pour un espace
annelé quelconque.}

\emph{Le cas visé.} Ces conditions sont satisfaites si $X$ est un polycylindre
compact de $\mathbb{C}^{n}$, muni des germes de fonctions holomorphes. En marge,
entouré : « premier espoir ».\footnote{« Espace de Stein » est d'abord écrit,
puis biffé et remplacé par « polycylindre analytique (compact) ». La
correction est juste : pour un compact de Stein quelconque, l'anneau
$\Gamma(K, \mathcal{O})$ n'est pas noethérien en général. Le théorème
correspondant est dû à J.~Frisch (\emph{Invent.\ Math.} 4, 1967) pour les
compacts semi-analytiques de Stein, et, à notre connaissance, Y.-T.~Siu (1969) a caractérisé les
compacts de Stein $K$ pour lesquels $\Gamma(K, \mathcal{O})$ est noethérien :
ce sont ceux dont l'intersection avec tout sous-ensemble analytique défini au
voisinage n'a qu'un nombre fini de composantes connexes. Les notes sont datées
d'après 1961 par l'inventaire et ne citent personne ; rien ne dit si elles
précèdent ces travaux.}

\section*{Les feuillets dactylographiés}

Sept feuillets de tapuscrits sont intercalés entre les notes, et n'ont rien à
voir avec elles. Six appartiennent à un même chapitre sur la dimension et la
codimension, dans les espaces topologiques puis dans les préschémas localement
noethériens ; ils sont lus ici dans l'ordre de leur pagination propre. Le
septième est d'un autre tapuscrit. On donne les énoncés et ce qu'il y a
ajouté.\footnote{La matière de ces feuillets figure, remaniée, dans le chapitre
0 d'EGA~IV, §14 (dimension combinatoire), et dans EGA~IV, §§4--6, à notre
connaissance ; on ne prétend pas identifier de quelle rédaction ils sont un
état.}

\pagerange{10}{10}
\subsection*{Dimension des espaces topologiques (page 10, tapuscrit p.~101)}

\textbf{Proposition 1.} (i) Si $Y$ est une partie d'un espace topologique $X$,
$\dim(Y) \leq \dim(X)$. (ii) Si $X$ est réunion finie de parties fermées
$X_{i}$, $\dim(X) = \sup_{i} \dim(X_{i})$. \textbf{Corollaire.} Si un voisinage
$U$ de $x$ est réunion de parties fermées $Y_{i}$ contenant $x$,
$\dim_{x} X = \sup_{i} \dim_{x}(Y_{i})$.

\textbf{Proposition 2.} Soit $\Phi$ l'ensemble des parties fermées
irréductibles de $X$. (i) $\dim(X) = \sup_{Y \in \Phi} \mathrm{codim}_{X}(Y)$.
(ii) $\dim(X) = \sup_{x} \dim_{x}(X)$.\footnote{Une assertion (iii) de la
proposition 1 est biffée à la machine. Le tapuscrit donne l'exemple d'une
application continue qui augmente la dimension : un espace discret à deux
points sur l'espace de Sierpiński.}

\pagerange{12}{12}
\subsection*{Espaces noethériens (page 12, tapuscrit p.~102)}

Soit $X$ noethérien et $T_{0}$. \textbf{Corollaire 2.} $\dim(X)$ est le
$\sup$ des $\dim_{x}(X)$ sur les points fermés. \textbf{Corollaire 3.} Pour
$Y$ fermé, $\dim(Y) + \mathrm{codim}_{X}(Y) \leq \dim(X)$.
\textbf{Corollaire 4.} Pour $Y \subset Z \subset T$ fermés irréductibles,
$\mathrm{codim}_{T}(Y) \geq \mathrm{codim}_{Z}(Y) + \mathrm{codim}_{T}(Z)$.
\textbf{Proposition 3.} $\dim(X) = 0$ si et seulement si $X$ est fini et
discret ; une partie fermée $Y$ est de codimension nulle si et seulement si
elle contient une composante irréductible de $X$.

\pagerange{14}{14}
\subsection*{La condition des chaînes (page 14, tapuscrit p.~104)}

C'est le feuillet où il est le plus intervenu. On dit que $X$ est
\emph{équidimensionnel} si ses composantes irréductibles ont toutes la même
dimension, \emph{équicodimensionnel} si ses points fermés ont tous la même
codimension, et qu'il satisfait la \emph{condition des chaînes} --- il est
caténaire --- si, pour $Y \subset Z$ fermés irréductibles, toutes les chaînes
saturées d'extrémités $Y$ et $Z$ ont même longueur.

\textbf{Proposition 5.} Soit $X$ un espace noethérien, $T_{0}$, de dimension
finie. Les conditions suivantes sont équivalentes :
\begin{enumerate}
\item[a)] deux chaînes maximales de parties fermées irréductibles de $X$ ont
même longueur ;
\item[c)] $X$ est équidimensionnel, et pour $Y \subset Z$ fermés
irréductibles, $\dim(Z) = \dim(Y) + \mathrm{codim}_{Z}(Y)$ ;
\item[d)] $X$ est équicodimensionnel, et pour $Y \subset Z$ fermés
irréductibles, $\mathrm{codim}_{X}(Y) = \mathrm{codim}_{Z}(Y) +
\mathrm{codim}_{X}(Z)$.
\end{enumerate}
Elles entraînent que $X$ est équidimensionnel, équicodimensionnel et
caténaire, mais cette dernière conjonction ne les entraîne pas. Sous ces
conditions, $\dim(Y) + \mathrm{codim}_{X}(Y) = \dim(X)$ pour toute partie
fermée irréductible $Y$ --- formule qu'il entoure --- et en particulier
$\dim(X) = \mathrm{codim}_{X}\{x\}$ pour tout point fermé $x$, comme le dit un
corollaire écrit de biais dans la marge.\footnote{Le tapuscrit donne quatre
conditions a), b), c), d), où b) et d) portaient « les points fermés de $X$
ont même codimension ». Il biffe ces mots dans les deux conditions et écrit
au-dessus « $X$ est équidimensionnel », avec en b) un renvoi marginal « et
équidimensionnel ». Les deux notions ne sont pas interchangeables, et aucune
forme de b) n'est équivalente à a). Pour c), l'équidimensionalité est la bonne
hypothèse, et c'est ce que le tapuscrit avait déjà. Pour d), la substitution
est fausse : si $R$ est un anneau de valuation discrète d'uniformisante $\pi$,
$X = \mathrm{Spec}\, R[x]$ est irréductible et caténaire, et vérifie la formule
de d), mais les chaînes maximales $(0) \subset (\pi x - 1)$ et
$(0) \subset (\pi) \subset (\pi, x)$ n'ont pas la même longueur ; on rétablit
donc l'équicodimensionalité du tapuscrit. Pour b), ni la forme tapée
(équicodimensionnel et caténaire) ni la forme corrigée (équidimensionnel et
caténaire) ne suffit --- le germe d'un plan et d'une droite concourants
réfute la première, $\mathrm{Spec}\, R[x]$ la seconde --- et même les deux
ensemble ne suffisent pas : sur l'espace fini à six points ordonné par
$P < w_{1}$, $Q < q' < w_{1}$, $P < p' < w_{2}$, les fermés étant les parties
stables par passage aux points plus petits, les deux composantes sont de
dimension 2, les deux points fermés de codimension 2, la condition des chaînes
est satisfaite, et pourtant $P \subset \overline{\{w_{1}\}}$ est une chaîne
maximale de longueur 1. On garde donc a), c), d), avec les lettres du
tapuscrit. Les deux derniers membres de l'égalité du corollaire marginal sont
lus sur les lettres et ne s'accordent pas avec ce qu'on attendrait ; on ne les
reprend pas.}

\pagerange{6}{6}
\subsection*{Préschémas algébriques sur un corps (page 6, tapuscrit p.~109)}

\textbf{Corollaire 2.} Soit $X$ un préschéma algébrique sur un corps $k$ et
$x \in X$. Alors $\dim_{x}(X) = \sup_{i} \dim(X_{i})$, les $X_{i}$ parcourant
les composantes irréductibles contenant $x$, et
\[
  \dim_{x}(X) = \dim \mathcal{O}_{X,x} + \mathrm{deg.tr}_{k}\, k(x).
\]
La seconde formule est ajoutée de sa main, là où le tapuscrit ne l'énonçait
que pour $x$ fermé, et elle est juste en général.\footnote{Elle se déduit de
la première, parce que chaque composante $X_{i}$ est équidimensionnelle et
caténaire : $\dim X_{i} = \mathrm{codim}_{X_{i}}\overline{\{x\}} +
\dim \overline{\{x\}}$, et $\dim \overline{\{x\}}$ est le degré de
transcendance de $k(x)$. Il écrit « $\dim k(x)/k$ ». La même formule revient
en marge, en regard d'un contre-exemple où $X = \mathrm{Spec}(K) \to
\mathrm{Spec}(B)$, $B$ de valuation discrète, montre qu'elle ne vaut pas hors
du cas algébrique.}

\textbf{Corollaire 3.} Pour un $k$-morphisme $f : X \to Y$ de préschémas
algébriques, $\dim \overline{f(X)} \leq \dim X$.

\pagerange{4}{4}
\subsection*{Produits, et l'inégalité des dimensions (page 4, tapuscrit p.~110)}

Le feuillet finit la démonstration de $\dim(X \times_{k} Y) = \dim(X) +
\dim(Y)$ pour des préschémas de type fini sur un corps, et écrit « en
particulier, $\dim(X \otimes_{k} K) = \dim(X)$ pour toute extension $K$ de
$k$ ». Il souligne « en particulier » d'un trait ondulé et écrit en marge :
« abusif ».\footnote{La remarque est juste : $K$ n'est pas de type fini sur
$k$ en général, et $\mathrm{Spec}(K)$ n'est pas dans le champ de la formule
qui précède ; l'égalité est vraie, mais ne s'en déduit pas « en particulier ».
Une note en biais, en regard du lemme de normalisation, parle des extensions
$K(z)$, $K(x)$, $K(y)$ et ne se lit qu'en fragments.}

\textbf{Proposition 8.} Soient $X$, $Y$ localement noethériens, $f : X \to Y$ un
morphisme, $y \in Y$ et $x$ le point générique d'une composante irréductible
de $f^{-1}(y)$. Alors $\mathrm{codim}_{X}\overline{\{x\}} \leq
\mathrm{codim}_{Y}\overline{\{y\}}$, c'est-à-dire
$\dim \mathcal{O}_{X,x} \leq \dim \mathcal{O}_{Y,y}$.

Il biffe l'hypothèse « de type fini » du tapuscrit, qui est en effet
inutile, et note en marge de simplifier la preuve par l'inégalité
\[
  \dim \mathcal{O}_{X,x} \leq \dim \mathcal{O}_{Y,y} + \dim \mathcal{O}_{f^{-1}(y),x},
\]
dont la proposition est le cas où le dernier terme est nul, puis :
« donner le cas d'égalité par \emph{platitude} ».\footnote{La marge écrit le
dernier terme $\dim \mathcal{O}_{x}$, ce qui rendrait l'inégalité triviale ;
le terme qui porte l'argument est la dimension de l'anneau local de la fibre,
et on le rétablit. L'inégalité vaut pour tout homomorphisme local d'anneaux
locaux noethériens, avec égalité quand il est plat : c'est exactement ce que
la marge demande d'ajouter.}

\pagerange{2}{2}
\subsection*{Codimension et images réciproques (page 2, tapuscrit p.~111)}

\textbf{Corollaire 1.} Soient $X$, $Y$ localement noethériens, $f : X \to Y$,
$Y'$ fermé irréductible dans $Y$, et $X'$ une composante irréductible de
$f^{-1}(Y')$ dont l'image est dense dans $Y'$. Alors
$\mathrm{codim}_{X}(X') \leq \mathrm{codim}_{Y}(Y')$.

\textbf{Corollaire 2.} Soit $X$ localement noethérien et $f$ une section de
$\mathcal{O}_{X}$ qui n'est identiquement nulle sur aucune composante
irréductible. Toute composante irréductible du lieu $Z_{f}$ des zéros de $f$
est de codimension 1.

Les deux « localement noethérien » sont de sa main. En regard de la fin de la
démonstration précédente, qui établit laborieusement une inégalité de rangs
d'idéaux premiers, il écrit « propositions laborieuses inutiles », et en
regard de la démonstration du corollaire 2, qui passe par le
\emph{Hauptidealsatz} de Krull, « affirmations inutiles, si on remarque que »,
la suite étant illisible.\footnote{Le premier corollaire est en effet la
proposition 8 appliquée au point générique de $X'$, qui est générique dans une
composante de la fibre au-dessus du point générique de $Y'$ ; c'est sans doute
ce qu'il juge inutile de démontrer à nouveau, mais la note ne le dit pas.}

\pagerange{8}{8}
\subsection*{Cycles premiers associés (page 8, un autre tapuscrit)}

Un feuillet paginé « 8 », frappé sur une machine sans accents, sur les points
associés d'un Module quasi-cohérent : $\mathrm{Ass} \coprod F_{i} = \bigcup_{i}
\mathrm{Ass}\, F_{i}$ ;

\textbf{Proposition 6.7.} Sur $X$ localement noethérien, une section
$f \in \Gamma(\mathcal{O}_{X})$ définit une homothétie injective de $F$ si et
seulement si $V(f) \cap \mathrm{Ass}\, F = \emptyset$ ;

\textbf{Corollaire 6.8.} Un Idéal cohérent $I$ est non diviseur de zéro dans un
Module cohérent $F$ si et seulement si $V(I) \cap \mathrm{Ass}\, F = \emptyset$ ;

\textbf{Proposition 6.9.} Soit $U$ un ouvert dense de $X$ localement
noethérien. $X$ est réduit si et seulement si $U$ l'est et $X$ n'a pas de
cycles premiers immergés ; ou encore si et seulement si $X$ n'a pas de cycles
premiers immergés et $\mathcal{O}_{X,x}$ est de longueur 1 aux points
génériques des composantes.

En marge de la démonstration de 6.7, qui renvoie au cas affine « où c'est dans
Bourbaki » : « démontrer directement ».\footnote{Le tapuscrit renvoie à
« IV 2.8.18 et 2.8.19 » et observe lui-même qu'« il est peut-être préférable
d'essayer de donner de telles démonstrations “intuitives”, plutôt que de se
réduire toujours à l'énoncé connu dans le cas affine » ; la note manuscrite va
dans le même sens.}

\pagerange{16}{17}
\section*{Un tore complexe qui n'est pas une variété abélienne (pages 16 et 17)}

\textbf{Énoncé.} Il existe un tore complexe $T$ de dimension 2, contenant un
sous-tore $T_{0}$ de dimension 1, qui n'est pas une variété abélienne. Comme
$T \to T/T_{0}$ est un fibré analytique localement trivial de fibre $T_{0}$ au
dessus d'un tore de dimension 1, et que les tores de dimension 1 sont des
courbes elliptiques, il s'ensuit qu'un fibré analytique dont la base et la
fibre sont des variétés de Hodge n'est pas nécessairement une variété de
Hodge.\footnote{Variété de Hodge : variété kählérienne compacte munie d'une
forme de Kähler à classe entière, c'est-à-dire, par le théorème de plongement
de Kodaira, variété projective. Ce sont les mots de la page.}

\textbf{Exemple.} Soit $k$ un réel irrationnel, et $\Gamma \subset
\mathbb{C}^{2}$ le sous-groupe engendré par
\[
  e_{1}, \quad ie_{1}, \quad e_{2}, \quad ke_{1} + ie_{2} .
\]
Ces quatre vecteurs sont $\mathbb{R}$-linéairement indépendants, $\Gamma$ est un
réseau, et $T = \mathbb{C}^{2}/\Gamma$ est un tore. La droite $V =
\mathbb{C}e_{1}$ rencontre $\Gamma$ en $\mathbb{Z}e_{1} + \mathbb{Z}ie_{1}$, et
$T_{0} = V/(\Gamma \cap V) = \mathbb{C}/\mathbb{Z}[i]$ est un sous-tore.\footnote{La
page écrit le quatrième générateur $e_{1} + ike_{2}$. Avec ce choix, $\Gamma$
contient $ike_{2} = (e_{1} + ike_{2}) - e_{1}$, donc
$\Gamma = (\mathbb{Z} + i\mathbb{Z})e_{1} \oplus (\mathbb{Z} + ik\mathbb{Z})e_{2}$,
et $T = \mathbb{C}/\mathbb{Z}[i] \times \mathbb{C}/(\mathbb{Z} + ik\mathbb{Z})$
est un produit de deux courbes elliptiques --- une variété abélienne. Le
raisonnement de la page échoue précisément là : si $x = \gamma e_{2} +
\delta(ike_{2})$ est pris sur l'axe $\mathbb{C}e_{2}$, ses deux premières
composantes sont nulles et n'entraînent pas la rationalité de $c'$ et $c''$ ;
et en effet $x = e_{2}$, $c = ik$ donnent $cx \in \Gamma$. En plaçant $k$ sur
l'autre coordonnée, comme ci-dessus, le calcul de la page passe, les rôles des
deux coordonnées étant échangés. C'est la correction que nous proposons ; le
fragment ne la contient pas.}

\emph{Démonstration.} Si $T$ était une variété abélienne, le théorème de
réductibilité de Poincaré fournirait un sous-tore complémentaire de $T_{0}$,
c'est-à-dire une droite complexe $V'$ supplémentaire de $V$ telle que
$\Gamma \cap V'$ soit un réseau de $V'$.\footnote{La page écrit « tel que $V$
soit $\Gamma$-propre, i.e.\ $V$ engendré par $\Gamma \cap V$ » ; c'est de $V'$
qu'il s'agit, comme le montre la suite. Le mot « propre » est lu avec doute.}
Il existerait alors $x \in \Gamma \cap V'$ non nul et un nombre complexe non
réel $c = c' + ic''$ tels que $cx \in \Gamma$ : deux vecteurs
$\mathbb{R}$-indépendants d'un réseau de la droite complexe $V'$.

Écrivons $x = \alpha e_{1} + \beta ie_{1} + \gamma e_{2} + \delta(ke_{1} +
ie_{2})$ avec $\alpha, \beta, \gamma, \delta \in \mathbb{Z}$ : ses coordonnées
sont
\[
  z_{1}(x) = (\alpha + k\delta) + i\beta, \qquad z_{2}(x) = \gamma + i\delta,
\]
et $(\gamma, \delta) \neq (0, 0)$ puisque $x \notin V$. Les coordonnées de
$cx$ sont
\[
  z_{1}(cx) = c'(\alpha + k\delta) - c''\beta + i\bigl(c''(\alpha + k\delta)
  + c'\beta\bigr),
  \qquad
  z_{2}(cx) = c'\gamma - c''\delta + i(c''\gamma + c'\delta),
\]
et, $cx$ étant dans $\Gamma$, elles doivent s'écrire $(\alpha' + k\delta') +
i\beta'$ et $\gamma' + i\delta'$ avec des entiers.

\begin{enumerate}
\item[1.] $c'\gamma - c''\delta$ et $c''\gamma + c'\delta$ sont entiers, et
$\gamma^{2} + \delta^{2} \neq 0$ : donc $c'$ et $c''$ sont rationnels.
\item[2.] $c''(\alpha + k\delta) + c'\beta$ est entier : donc $c''k\delta$ est
rationnel, et, $k$ étant irrationnel et $c'' \neq 0$, $\delta = 0$.
\item[3.] $c'\alpha - c''\beta = \alpha' + k\delta'$ : donc $k\delta'$ est
rationnel, $\delta' = 0$, et $c''\gamma + c'\delta = c''\gamma = 0$, d'où
$\gamma = 0$.
\end{enumerate}

Mais alors $\gamma = \delta = 0$, et $x$ est dans l'espace engendré par
$e_{1}$ et $ie_{1}$, c'est-à-dire dans $V$ : absurde.\footnote{C'est, pas pour
pas, la suite d'arguments de la page 17 --- rationalité par deux composantes
entières, annulation de $\delta$ par l'irrationalité de $k$, puis de $\gamma$,
contradiction « $x$ serait dans l'espace engendré par $e_{1}$ et $ie_{1}$ » ---
appliquée au réseau corrigé. Plus généralement, pour le réseau engendré par
$e_{1}, ie_{1}, e_{2}, be_{1} + ie_{2}$, $b \in \mathbb{C}$, le même calcul
montre que $T$ est abélien si et seulement si $b \in \mathbb{Q}(i)$. Les tores
de dimension 2 ainsi obtenus sont de dimension algébrique 1 ; le nom n'est pas
sur la page.}

\end{document}
